\documentclass[10pt,letterpaper]{article}

\usepackage[T1]{fontenc}
\usepackage[utf8]{inputenc}
\usepackage{lmodern}
\usepackage[margin=0.9in,top=0.8in]{geometry}
\usepackage{enumitem}
\usepackage{tabularx}
\usepackage{titlesec}
\usepackage[hidelinks]{hyperref}
\usepackage{xcolor}

\hypersetup{colorlinks=true,urlcolor=blue,linkcolor=blue}
\pagestyle{empty}
\setlength{\parindent}{0pt}

% Section headings: bold caps with a rule underneath
\titleformat{\section}{\large\bfseries\uppercase}{}{0pt}{}[\vspace{-0.6em}\rule{\linewidth}{0.5pt}]
\titlespacing*{\section}{0pt}{0.9em}{0.5em}

% Content blocks are indented under each heading
\newenvironment{block}{\begin{list}{}{\leftmargin=1.5em\rightmargin=0pt\itemsep=0pt\topsep=0pt}\item[]}{\end{list}}

% Dot bullets, as in the template
\newenvironment{bullets}{\begin{itemize}[label=$\cdot$,leftmargin=0.9em,itemsep=0pt,parsep=0pt,topsep=-0.3em]}{\end{itemize}\vspace{0.55em}}

% \entry{Org}{Dates}{Role}{Location}
\newcommand{\entry}[4]{%
  \begin{tabularx}{\linewidth}{@{}>{\raggedright\arraybackslash}X r@{}}
    \textbf{#1} & #2 \\
    \textit{#3} & \textit{#4}
  \end{tabularx}\par\vspace{0.3em}}
\newcommand{\me}{\textbf{Zhian Ruan}}
\newcommand{\paper}[1]{\href{#1}{[Paper]}}

\begin{document}

\begin{center}
  {\bfseries\fontsize{24}{28}\selectfont ZHIAN (ARIA) RUAN}\\[0.6em]
  \href{https://ariar1225.github.io/}{https://ariar1225.github.io/} $|$ ariaruan@stanford.edu $|$
  \href{https://github.com/ariaR1225}{GitHub} $|$
  \href{https://www.linkedin.com/in/aria-ruan-11055a201/}{LinkedIn}
\end{center}

\section{Education}
\begin{block}
\entry{Stanford University}{09/2026-now}{M.S. in Electrical Engineering, Advisor: Prof. Shuran Song}{}
\entry{Northwestern University}{09/2022-06/2025}{B.S. in Mechanical Engineering, Concentration in Robotics, Magna Cum Laude}{GPA 3.947/4.0}
\entry{University of Illinois Urbana-Champaign}{08/2020-05/2022}{B.S. in Mechanical Engineering, transferred to Northwestern University}{GPA 4.00/4.00}
\end{block}

\section{Research Interest}
\begin{block}
My research lies at the intersection of robot learning and system integration. I focus on synergizing
hardware--software--learning co-design, multimodal model architectures, and reinforcement learning to
enable robust, adaptive, and generalizable robot policies, tightly integrating sensing, learning, and
control into efficient, reliable, ready-to-deploy robotic systems for the real world.
\end{block}

\section{Publications}
\begin{enumerate}[leftmargin=2.6em,itemsep=0.5em,topsep=0pt]
  \item \textit{Whole Body Control.} Under review.
  \item Anthony Goeckner, Yao Liu, \me, Alessandro Pinto, Qi Zhu, \textit{Quiet Multi-Agent Systems via
        Predictive World Modeling and Selective Communications.} Under review.
  \item Simon Sinong Zhan, Philip Wang, Yao Liu, Yiyan Peng, Yiqi Lyu, Zinan Wang, Qineng Wang,
        \me, Xiangyu Shi, Xinyu Cao, Frank Yang, Zhenyang Ni, Kangrui Wang, Ruohan Zhang, Huajie Shao,
        Manling Li, Qi Zhu, \textit{SENTINEL: A Multi-Level Formal Framework for Safety Evaluation of
        Foundation Model-based Embodied Agents.} Neural Information Processing Systems, Evaluations \&
        Datasets Track (NeurIPS 2026). \paper{https://arxiv.org/abs/2510.12985}
  \item \me, Yizhi Zhou, \textit{Distributed Invariant Unscented Kalman Filter based on Inverse
        Covariance Intersection with Intermittent Measurements.} American Control Conference (ACC 2026).
        \paper{https://arxiv.org/abs/2409.17997}
  \item Simon Sinong Zhan, Qingyuan Wu, \me, Frank Yang, Philip Wang, Yixuan Wang, Ruochen Jiao,
        Chao Huang, Qi Zhu, \textit{Inverse Delayed Reinforcement Learning.} arXiv preprint (2024).
        \paper{https://arxiv.org/abs/2412.02931}
\end{enumerate}

\section{Research Experience}
\begin{block}
\entry{Robotics and Embodied Artificial Intelligence Lab (REAL), Stanford University}{03/2026-now}{Advisor: Prof. Shuran Song}{Stanford, CA}
\begin{bullets}
  \item Force-sensing UMI data collection and diffusion policies for whole-body mobile manipulation
\end{bullets}
\entry{Design Automation of Intelligent Systems (IDEAS) Lab, Northwestern University}{06/2024-06/2025}{Advisor: Prof. Qi Zhu}{Evanston, IL}
\begin{bullets}
  \item LLM embodied-agent safety evaluation, inverse delayed RL, and quiet multi-agent planning
\end{bullets}
\entry{Neurorehabilitation and Neural Engineering Lab, Shirley Ryan AbilityLab}{06/2023-08/2023}{Advisor: Dr. Jos\'e L. Pons}{Chicago, IL}
\begin{bullets}
  \item Transparency control for a lower-limb exoskeleton, later part of a \href{https://www.science.org/doi/10.1126/scirobotics.adz9628}{Science Robotics paper}
\end{bullets}
\entry{Center for Robotics and Biosystems, Northwestern University}{03/2023-08/2023}{Advisor: Prof. Malcolm MacIver}{Evanston, IL}
\begin{bullets}
  \item 4\,m/s embedded line-following robot with PID control and a Wi-Fi command link
\end{bullets}
\entry{Center for Bionics, Shirley Ryan AbilityLab}{08/2022-06/2023}{Advisor: Dr. R. James Cotton}{Chicago, IL}
\begin{bullets}
  \item Gait analysis from EMG muscle synergies via non-negative matrix factorization
\end{bullets}
\entry{Human Dynamics and Controls Lab, University of Illinois Urbana-Champaign}{07/2021-07/2022}{Advisor: Prof. Elizabeth T. Hsiao-Wecksler}{Urbana, IL}
\begin{bullets}
  \item EMG database and quantitative measures of abnormal muscle behavior in MATLAB
\end{bullets}
\end{block}

\section{Professional Experience}
\begin{block}
\entry{RoboForce}{10/2025-now}{Member of Technical Staff, Robotics}{Milpitas, CA}
\begin{bullets}
  \item Led end-of-arm tooling (EOAT) and data collection system design for model training, reducing camera
        latency and enabling cross-platform time synchronization
  \item Validated and brought up perception and localization sensors, including cameras, LiDAR, motion capture,
        and VIO, across lighting and terrain conditions
  \item Selected and tested joint motors, and led robot bringup with ROS 2 device drivers in C++ over RS485,
        CANopen, and EtherCAT (CiA 402) with a thread-safe, real-time communication stack
  \item Built a force-torque sensing (FTS) embedding for a $\pi_{0.5}$ torquing policy: a causal TCN/GRU encoder over
        2\,s of native 200\,Hz wrench, aligned across asynchronous sensor, camera, and action clocks, compressing contact
        into a bounded 6-D embedding that reaches 0.90 macro recall on contact phase (+0.17 over a no-memory baseline)
        across 1,900+ sessions
\end{bullets}

\entry{Tesla, Energy R\&D Test Infrastructure}{07/2025-10/2025}{Test Hardware Controls Engineer}{Palo Alto, CA}
\begin{bullets}
  \item Engineered a centralized Python resource manager for test infrastructure, automating testing across 3
        energy product lines with multi-site coordination
  \item Designed a full-stack system with a REST API frontend for remote control of test setups and a real-time
        PLC backend over Modbus, unifying manual and automated workflows
\end{bullets}

\textbf{Internships:} Ginkgo Bioworks, system integration \& computer vision (06/2024-09/2024);
Tesla Power Electronics, software automation (01/2024-05/2024)
\end{block}

\section{Selected Projects}
\begin{block}
\entry{Robotic Dexterity Hand, Embedded Control Software}{01/2025-06/2025}{Advisor: Prof. Kevin Lynch, HAND ERC}{}
\begin{bullets}
  \item Led control software for an 8-DOF tendon-driven hand (20-student team): SPI Hall-effect sensing and CAN
        across 9 MCUs, enabling 10\,kHz field-oriented BLDC control with space-vector PWM
  \item Built motion-control libraries with forward/inverse kinematics and null-space optimization, reaching
        $<$2\,mm positioning accuracy, $<$0.05\,s force settling time, and 10\,Hz position bandwidth on series-elastic actuators
  \item Created a MediaPipe hand-tracking teleoperation GUI with ROS 2 and a URDF digital twin in RViz
\end{bullets}


\end{block}

\section{Awards}
\begin{block}
Magna Cum Laude, Northwestern University \hfill 06/2025\\
McCormick School of Engineering High Honors, Northwestern University \hfill 2023-2024\\
McCormick School of Engineering Honors, Northwestern University \hfill 2022-2023\\
Durabird Eastern Collegiate Badminton Championships, Division 2B Winner \hfill 2023\\
Edmund J. James Scholar \& Dean's List, Grainger College of Engineering, UIUC \hfill 2020-2022\\
Tau Beta Pi, National Engineering Honor Society \hfill 2021\\
Alpha Lambda Delta, National Honor Society of First-Year Success \hfill 2021
\end{block}

\section{Services}
\begin{block}
Reviewer of CoRL 2026 Workshop on \href{https://scaling-h2r-corl.github.io/}{Scaling H2R}\par\vspace{0.4em}
Reviewer of IEEE RO-MAN 2026\par\vspace{0.4em}
Grader of GEN\_ENG 205-3: Engineering Analysis 3 (Modeling and Dynamics), Northwestern (Spring 2025)\par\vspace{0.4em}
Course assistant of CS 111: Fundamentals of Computer Programming I, Northwestern (Fall 2023)\par\vspace{0.4em}
Student consultant for Chicago Field Studies, Student Internship Program (Fall 2023)\par\vspace{0.4em}
Executive of the Northwestern Badminton Club, 300+ members (2023)\par\vspace{0.4em}
Executive board of the Tau Beta Pi District Conference, 30+ chapters (2022)\par\vspace{0.4em}
Mentor for FRC \href{https://www.thebluealliance.com/team/6487}{Team 6487} (2020-now)
\end{block}

\section{Technical Skills}
\begin{block}
\begin{tabularx}{\linewidth}{@{}l >{\raggedright\arraybackslash}X@{}}
\textbf{Computer Languages} & Python, C/C++, Java, MATLAB/Simulink, SQL \\
\textbf{Robotics} & ROS 2, MuJoCo, IsaacGym, Gazebo, RViz, forward/inverse kinematics, SLAM, PID \\
\textbf{Tools \& Libraries} & PyTorch, TensorFlow, NumPy, SciPy, OpenCV, Git, Docker, Linux \\
\textbf{Hardware} & CAD, 3D printing, PCB design, PLC, microcontrollers (ARM Cortex-M4/M7), CAN, EtherCAT, RS485, SPI, I2C, UART \\
\end{tabularx}
\end{block}

\end{document}
